\documentstyle[11pt,fleqn]{article}
\textheight=220mm
\textwidth=160mm
\voffset=-16mm
\hoffset=-16mm
\parskip=4pt plus2pt minus2pt
\begin{document}

\begin{enumerate}

\item  Given the M/M/2/2 system transitions:

\begin{figure}[hbt]\centering
\setlength{\unitlength}{0.16mm}
\begin{picture}(720,120)(0, 60)

\thicklines

\put(120,  120){\oval(60, 60)} 
\put(120,  120){\makebox(0, 0){\sf 0}}
\put(150,  140){\vector(1, 0){60}} 
\put(180,  150){\makebox(0, 0){$\lambda$}}


\put(240,  120){\oval(60, 60)} 
\put(240,  120){\makebox(0, 0){\sf 1}}
\put(270,  140){\vector(1, 0){60}} 
\put(300,  150){\makebox(0, 0){$\lambda$}}
\put(210,  100){\vector(-1, 0){60}} 
\put(180,   80){\makebox(0, 0){$\mu$}}

\put(360,  120){\oval(60, 60)} 
\put(360,  120){\makebox(0, 0){\sf 2}}
\put(330,  100){\vector(-1, 0){60}} 
\put(300,   80){\makebox(0, 0){2*$\mu$}}

\end{picture}
\label{M/M/1/K}
\end{figure}

\begin{enumerate}

\item show that P$_0$, the probability of no customers in the system =
1/(1 + 2$\rho$ + 2$\rho^2$), where {\bf $\rho$ = $\lambda$/2$\mu$},
and derive formulae for P$_1$ and P$_2$.   

\item show that N, the average number of customers in the
system, is given by N =  2$\rho$$\frac{1+2\rho}{1 + 2\rho + 2\rho^2}$

\item derive a formula for T, the average customer through time 

\end{enumerate}

\newpage


\item Given system of 4 M/M/1 nodes with:

\begin{enumerate}

\item Service rates:

Node 1 completes 10 jobs/second;
Node 2 completes 10 jobs/second;
Node 3 completes 20 jobs/second;
Node 4 completes 20 jobs/second.

\item Inputs: 
$\lambda_1$ = 1; $\lambda_2$ = 2; 
$\lambda_3$ = 3; $\lambda_4$ = 4 --- all measured in jobs/second.

\item Recirculation:

>From node 1: jobs go back to node 1 with probability 0.4
and to node 2 with probability 0.2.

>From node 2: jobs go back to node 2 with probability 0.2
and to node 3 with probability 0.2.

>From node 3: jobs go back to node 3 with probability 0.1
and to node 4 with probability 0.2.

>From node 4: jobs go to 1 with probability 0.5.

\end{enumerate}

\noindent
{\bf For each node compute}

(i) the effective arrival rate\\
(ii) the utilisation\\
(iii) the average number in that node and
the average value queueing at that node.\\
(iv) the average time spent at that node and
the average time spent queueing at that node.\\
(v) the rate at which jobs leave the system at that node.\\

\noindent
{\bf For the whole system, compute}

(i) the arrival rate into the system as a whole\\
(ii) the expected number of jobs in the system\\
(iii) the expected time spent by a job in the system\\
(iv) the bottleneck device in the system\\

Run the model doubling the input rates.
Do the utilisations double? Do the through times double?
Which is the bottleneck device?
When does the system saturate
(i.e. contain a device with 100\% utilisation)?


\newpage

\item Consider an open Jackson network comprising
three M/M/1 systems, N$_1$, N$_2$, N$_3$, with mean service rates 

\begin{center}
$\mu_1$ = 4 \hspace*{0.5in} $\mu_2$ = 4 \hspace*{0.5in} $\mu_3$ = 6
\end{center}

and mean customer arrival rates of

\begin{center}
$\lambda_1$ = 0.75 \hspace*{0.5in} 
$\lambda_2$ = 0.50 \hspace*{0.5in} 
$\lambda_3$ = 0.75
\end{center}

both measured in customers per minute.

On completion of service at N$_1$,
50\% of the jobs queue again at N$_1$,
25\% of the jobs queue at N$_2$, and
25\% of the jobs queue at N$_3$.

On completion of service at N$_2$,
25\% of the jobs queue at N$_1$ and
75\% of the jobs queue at N$_3$.

On completion of service at N$_3$,
50\% of the jobs queue again at N$_3$
and 50\% of the jobs leave the system.

\begin{enumerate}

\item Show that the effective flow rates at the three nodes are:

\begin{center}
$\Lambda_1$ = 2 \hspace*{0.5in} 
$\Lambda_2$ = 1 \hspace*{0.5in} 
$\Lambda_3$ = 4
\end{center}

measured in customers per minute.

\item Hence find individual node utilisations 

\item ... the expected number of customers at each node, and 
the expected number of customers in the system as a whole,


\item ... and average time spent by a customer in the system.  

\end{enumerate}

\newpage

\item Consider the 4-node queueing described below.
Jobs arrive externally at nodes 1 and 2 with equal rates
$\lambda$/2. The routeing is as follows:

\begin{itemize}
\item from node 1, jobs go to node 3 with probability 0.3 
and to node 4 with probability 0.7

\item from node 2, jobs go to node 3 with probability 0.7 
and to node 4 with probability 0.3

\item from node 3, jobs go back to node 1 with probability 0.8 
and leave the network with probability 0.2 

\item from node 4, jobs go back to node 2 with probability 0.8 
and leave the network with probability 0.2 

\end{itemize}

\begin{enumerate}

\item Draw a picture of the network showing the four nodes
and the flows amongst them

\item Write down the flow rate equations for each node

\item Show that

\begin{center}
$\Lambda_1$ + $\Lambda_2$ =  $\Lambda_3$ + $\Lambda_4$ = 5$\times$$\lambda$
\end{center}

where $\Lambda_1$, $\Lambda_2$, $\Lambda_3$, and $\Lambda_4$ 
are the actual flow rates at the 4 nodes respectively


\item ... and hence that $\Lambda_1$ = $\Lambda_2$ =  $\Lambda_3$ = $\Lambda_4$


\item If the average service times are all equal, give 
and justify an upper bound for the $\Lambda_k$.

\end{enumerate}

\newpage

\item {\bf Buzen's algorithm for the closed central server model.}

Modify the first example in your notes by 
adding in a fourth device equivalent to device 3 
which shares its workload (i.e. P3  and P4 are both 0.1).

Modify this example again by 
adding in fourth/fifth/sixth devices
each equivalent to device 3 and which share its workload 
(i.e. P3, P4, P5, P6 are all 0.05).

\end{enumerate}

Any comments on the results?

\newpage

First two reminders of how to work out some well-known 
and elementary formulae used in queueing theory.

\subsection{1/1-x}

To find the sum of 1 + x + x$^2$ + x$^3$ + ....
Its not obvious so let's call it S. Then, by definition,

\begin{tabular}{r cc cc cc cc ccc}
S & = & 1 & + & x & + & x$^2$ & + & x$^3$ & + & x$^4$ & .... \\
\end{tabular}

\noindent
We notice that each term on the right is greater by x than the one
before.
Hmmh. This suggests we look at  S.x (S multiplied by x), 
more conveniently written the other way round as x.S.

\begin{tabular}{r cc cc cc cc ccc}
x.S & = & x & + & x$^2$ & + & x$^3$ & + & x$^4$ &  .... \\
\end{tabular}

\noindent
The right hand sides of S and x.S are almost the same
(its helpful to ``get the columns right'').


\begin{tabular}{r cc cc cc cc ccc}
S   & = & 1 & + & x & + & x$^2$ & + & x$^3$ & + & x$^4$ & .... \\
x.S & = &   &   & x & + & x$^2$ & + & x$^3$ & + & x$^4$ &  .... \\
\end{tabular}


\noindent
What if we subtract? 

\begin{tabular}{r cc cc cc cc ccc c}
S   & = & 1 & + & x & + & x$^2$ & + & x$^3$ & + & x$^4$ & .... \\
x.S & = &   &   & x & + & x$^2$ & + & x$^3$ & + & x$^4$ &  .... \\
\hline
S-x.S & = & 1 & \\
\hline
\end{tabular}

\noindent
Hence  S-x.S = 1, or S = 1/(1-x). So that we have shown that

\begin{table}[hbt]\centering
\begin{tabular}{||r cc cc cc cc ccc||}
\hline
&&&&&&&&&&&\\
\hspace*{0.25in}
1/(1-x)
  & = & 
1 & + & x & + & x$^2$ & + & x$^3$ & + & x$^4$ & .... 
\hspace*{0.25in}\\
&&&&&&&&&&&\\
\hline
\end{tabular}
\end{table}


\newpage



\subsection{1/(1-x)$^2$}

We also need to recognise the expansion of $\frac{1}{(1-x)^2}$.
We start off with the obvious fact that $\frac{1}{(1-x)^2}$
= $\frac{1}{(1-x)}$ $\times$ $\frac{1}{(1-x)}$, or S $\times$ S,
or as we have just shown 

\begin{center}
(1 + x + x$^2$ + x$^3$ + ....) $\times$ (1 + x + x$^2$ + x$^3$ + ....)
\end{center}

\noindent
the best way to proceed is to break the problem up into smaller steps


\begin{tabular}{r cc cc cc cc cc cc cc}
\multicolumn{10}{l}{$\frac{1}{(1-x)^2}$} \\
=  & 1 & $\times$ &(& 1 & + & x & + & x$^2$ & + & x$^3$ & + & x$^4$ &
  .... &)\\
+  & x & $\times$ &(& 1 & + & x & + & x$^2$ & + & x$^3$ & + & x$^4$ &
  .... &)\\
+  & x$^2$ & $\times$ &(& 1 & + & x & + & x$^2$ & + & x$^3$ & + &
  x$^4$ & .... &)\\
+  & x$^3$ & $\times$ &(& 1 & + & x & + & x$^2$ & + & x$^3$ & + & x$^4$ & .... &)\\
+  & x$^4$ & $\times$ &(& 1 & + & x & + & x$^2$ & + & x$^3$ & + &
  x$^4$ & .... &)\\
+ & \multicolumn{10}{l}{.........................................}\\
\end{tabular}

\noindent
Carrying through the multiplications line by line and taking care to
align so that we have columns of the same power, we get 

\begin{tabular}{r cc cc cc cc cc cc cc}
\multicolumn{10}{l}{$\frac{1}{(1-x)^2}$} \\
=  & 1 & + & x & + & x$^2$ & + & x$^3$ & + & x$^4$ &  .... &\\
   &   & + & x & + & x$^2$ & + & x$^3$ & + & x$^4$ &  .... &\\
   &   &   &   & + & x$^2$ & + & x$^3$ & + & x$^4$ &  .... &\\
   &   &   &   &   &       & + & x$^3$ & + & x$^4$ &  .... &\\
   &   &   &   &   &       &   &       & + & x$^4$ &  .... &\\
\multicolumn{10}{l}{.........................................}\\
\hline
= &  1 & + & 2.x & + & 3.x$^2$ & + & 4.x$^3$ & + & 5.x$^4$ &  .... &\\
\hline
\end{tabular}

\noindent
So that we have shown that

\begin{table}[hbt]\centering
\begin{tabular}{||r cc cc cc cc ccc||}
\hline
&&&&&&&&&&&\\
\hspace*{0.25in}
1/(1-x)$^2$
  & = & 
       1 & + & 2.x & + & 3.x$^2$ & + & 4.x$^3$ & + & 5.x$^4$ & .... 
\hspace*{0.25in}\\
&&&&&&&&&&&\\
\hline
\end{tabular}
\end{table}

\newpage
\section*{The exponential nature of random arrivals}

Throughout our look at systems of queues we are assuming that
our arrival and service patterns do not depend upon
the state of the system (i.e. they remain the same whether the system
is heavily loaded or lightly loaded).

We start off with the reasonable assumption that the number of
arrivals in a time slot (period of observation) 
depends upon the length of the time slot.
 
\begin{tabular}{l c l}
\rule{1in}{0.05in} & & expect n arrivals \\
\rule{2in}{0.05in} & & expect 2n arrivals \\
\rule{3in}{0.05in} & & expect 2n arrivals \\
\multicolumn{3}{c}{.............}\\
\end{tabular}

\noindent
If the arrival rate is $\lambda$ then we expect 
to average $\lambda$.t arrivals in a time slot of width t.

Now let's look at the probability of NO arrivals in a time slot 
of width t. We call this probability P(t).
So by definition the probability of no arrivals in a time slot 
of width t+$\delta$t is P(t+$\delta$t).

Now we get no arrivals in a time slot of width 
t+$\delta$t by having no arrival in the slot of width t
(and the probability of this happening is P(t))
and having no arrival in the next period of length $\delta$t.
And since the probability of an arrival in a period of length $\delta$t
is $\lambda$.$\delta$t (our starting assumption), 
then the probability of no arrival
in the period $\delta$t is (1- $\lambda$.$\delta$t).
So that 

\begin{center}
P(t+$\delta$t) = P(t) $\times$ (1 - $\lambda$.$\delta$t)
\end{center}

But the Taylor expansion of P(t+$\delta$t) is

\begin{center}
P(t+$\delta$t)  =  P(t)  +  $\delta$t.P'(t) + O($\delta$t$^2$) ....
\end{center}

Combining these last two formulae we get

\begin{center}
P(t)  +  $\delta$t.P'(t) .... =  P(t) - $\lambda$.$\delta$t.P(t)
\end{center}


As $\delta$t gets vanishingly small, we may neglect second order terms,
so 

\begin{center}
$\delta$t.P'(t) = - $\lambda$.$\delta$t.P(t)
\end{center}

i.e. P'(t) = - $\lambda$.P(t), the solution to which is 


\begin{table}[hbt]\centering
\begin{tabular}{|| cc c||}
\hline
&&\\
\hspace*{0.25in} P(t)   & = &  e$^{-\lambda t}$ \hspace*{0.25in}\\
&&\\
\hline
\end{tabular}
\end{table}

Remember that P(t) is is the probability of no arrivals in 
a time period of length t, so the the probability of 
some arrivals in this time period is (1 - e$^{-\lambda t}$).

\newpage
\begin{enumerate}

\item Given the M/M/1/2 system transitions:

\begin{figure}[hbt]\centering
\setlength{\unitlength}{0.16mm}
\begin{picture}(720,120)(0, 60)

\thicklines

\put(120,  120){\oval(60, 60)} 
\put(120,  120){\makebox(0, 0){\sf 0}}
\put(150,  140){\vector(1, 0){60}} 
\put(180,  150){\makebox(0, 0){$\lambda$}}


\put(240,  120){\oval(60, 60)} 
\put(240,  120){\makebox(0, 0){\sf 1}}
\put(270,  140){\vector(1, 0){60}} 
\put(300,  150){\makebox(0, 0){$\lambda$}}
\put(210,  100){\vector(-1, 0){60}} 
\put(180,   80){\makebox(0, 0){$\mu$}}

\put(360,  120){\oval(60, 60)} 
\put(360,  120){\makebox(0, 0){\sf 2}}
\put(330,  100){\vector(-1, 0){60}} 
\put(300,   80){\makebox(0, 0){$\mu$}}

\end{picture}
\label{M/M/1/2}
\end{figure}

\begin{enumerate}

\item show that P$_0$, the probability of no customers in the system =
1/(1 + $\rho$ + $\rho^2$), and derive formulae for P$_1$ and P$_2$. 

This system has only 3 states: S$_0$, S$_1$, and S$_2$.

P$_0$ + P$_1$ + P$_2$ = 1

The flow balance equations are 

\begin{tabular}{l c l ccc l c l}
$\lambda$.P$_1$ & = & $\mu$.P$_0$ 
   && or && P$_1$ & = & $\rho$.P$_0$\\
$\lambda$.P$_2$ & = & $\mu$.P$_1$ 
   && or && P$_2$ & = & $\rho^2$.P$_0$\\
\end{tabular}

Thus 

\begin{tabular}{l c l ccc l c l}
1  & = & P$_0$  & + & P$_1$ & + & P$_2$ \\
    & = & P$_0$  & + & $\rho$.P$_0$ & + & $\rho^2$.P$_0$ \\
    & = & \multicolumn{7}{l}{P$_0$(1 + $\rho$ + $\rho^2$)}\\[0.1in]
P$_0$ & = & \multicolumn{7}{l}{1/(1 + $\rho$ + $\rho^2$)}\\
P$_1$ & = & \multicolumn{7}{l}{$\rho$/(1 + $\rho$ + $\rho^2$)}\\
P$_2$ & = & \multicolumn{7}{l}{$\rho^2$/(1 + $\rho$ + $\rho^2$)}\\
\end{tabular}

Examples: if $\rho$ = 0.1, 0.5, and 0.9 then
(P$_0$, P$_1$, P$_2$) triples are (1/1.11, 0.1/1.11, 0.01/1.11),
(1/1.75,0.5/1.75,0.25/1.75),
and (1/2.71,0.9/2.71,0.81/2.71) respectively.
The chances of the system being full and an arrival rejected 
are 0.01/1.11, 0.25/1.75, and 0.81/2.71 respectively
(approx 1\%, 10\%, and 30\%).


\item show that N, the average number of customers in the
system, is given by N =  

\begin{tabular}{l c l ccc l c l}
N  & = & 0.P$_0$  & + & 1.P$_1$ & + & 2.P$_2$ \\
    & = & \multicolumn{7}{l}{$\rho$.P$_0$ + 2.$\rho^2$.P$_0$} \\
    & = & \multicolumn{7}{l}{$\rho$.(1+2.$\rho$).P$_0$}  \\
    & = & \multicolumn{7}{l}{$\rho$.(1+2.$\rho$)/(1 + $\rho$ + $\rho^2$)}\\[0.1in]
\end{tabular}

N$_q$ = 0.P$_0$ + 0.P$_1$ + 1.P$_2$ = P$_2$ = $\rho^2$/(1 + $\rho$ + $\rho^2$)


\item derive a formula for T, the average customer through time 

By Little's Law, T = N/$\lambda$', where the true arrival rate is
$\lambda$' = $\lambda$.(1-P$_2$) 
(with probability P$_2$ an arrival is discarded) 
= $\lambda$.(1 + $\rho$)/(1 + $\rho$  + $\rho^2$).

This is nice as the complicated terms cancel and we are left with:

T = E[s].(1+2.$\rho$)/(1+$\rho$)


\end{enumerate}

\newpage

\item Given the M/M/2/2 system transitions:

\begin{figure}[hbt]\centering
\setlength{\unitlength}{0.16mm}
\begin{picture}(720,120)(0, 60)

\thicklines

\put(120,  120){\oval(60, 60)} 
\put(120,  120){\makebox(0, 0){\sf 0}}
\put(150,  140){\vector(1, 0){60}} 
\put(180,  150){\makebox(0, 0){$\lambda$}}


\put(240,  120){\oval(60, 60)} 
\put(240,  120){\makebox(0, 0){\sf 1}}
\put(270,  140){\vector(1, 0){60}} 
\put(300,  150){\makebox(0, 0){$\lambda$}}
\put(210,  100){\vector(-1, 0){60}} 
\put(180,   80){\makebox(0, 0){$\mu$}}

\put(360,  120){\oval(60, 60)} 
\put(360,  120){\makebox(0, 0){\sf 2}}
\put(330,  100){\vector(-1, 0){60}} 
\put(300,   80){\makebox(0, 0){2*$\mu$}}

\end{picture}
\label{M/M/1/K}
\end{figure}

\begin{enumerate}

\item show that P$_0$, the probability of no customers in the system =
1/(1 + 2$\rho$ + 2$\rho^2$), where {\bf $\rho$ = $\lambda$/2$\mu$},
and derive formulae for P$_1$ and P$_2$.

\item show that N, the average number of customers in the
system, is given by N =  2$\rho$$\frac{1+2\rho}{1 + 2\rho + 2\rho^2}$


\item derive a formula for T, the average customer through time 

\end{enumerate}


\newpage

\item Consider four variants on the same theme. 

\begin{enumerate}

\item M/M/1 system (one queue, one server with
          mean service rate of 1 job/minute)

\item an M/M/1/2 system (at most 2 jobs in the system,
so the maximum queue length is 1,
one server with mean service rate of 1 job/minute)

\item M/M/2 system (one queue, two servers each with a
          mean service rate of 0.5 jobs/minute)

\item M/M/2/2 system (one queue of maximum capacity 1, 
    two servers each with a mean service rate of 0.5 jobs/minute)

\end{enumerate}

For arrival rates $\lambda$ of 0.1, 0.5, and 0.9 jobs per minute,
tabulate $\rho$, N, and T (the utilisation, 
expected number in the system, and
expected through times).
Any conclusions you can draw?

Reminder:

\begin{tabular}{lc lc lc l}
{\bf MM1} & & $\rho$ = ${\lambda}$/${\mu}$ &&
          N  = $\rho$/(1-$\rho$)&&
          T =  E[s]/(1-$\rho$)\\[0.1in]

{\bf MM1/2} & & $\rho$ = ${\lambda}$/${\mu}$ &&
          N  = $\rho$(1+2$\rho$)/(1+$\rho$+$\rho^2$)&&
          T =  E[s](1+2$\rho$)/(1+$\rho$)\\[0.1in]

{\bf MM2} & & $\rho$ = ${\lambda}$/2${\mu}$ &&
          N  = 2$\rho$/(1-$\rho^2$)&&
          T =  E[s]/(1-$\rho^2$)\\[0.1in]

{\bf MM2/2} && $\rho$ = $\lambda$/2$\mu$&&
N     =  2$\rho$(1+2$\rho$)/(1+2$\rho$+2$\rho^2$) &&
T     =  E[s] \\[0.1in]
\end{tabular}

\newpage


\item Consider an open Jackson network comprising
three M/M/1 systems, N$_1$, N$_2$, N$_3$, with mean service rates 

\begin{center}
$\mu_1$ = 4 \hspace*{0.5in} $\mu_2$ = 4 \hspace*{0.5in} $\mu_3$ = 6
\end{center}

and mean customer arrival rates of

\begin{center}
$\lambda_1$ = 0.75 \hspace*{0.5in} 
$\lambda_2$ = 0.50 \hspace*{0.5in} 
$\lambda_3$ = 0.75
\end{center}

both measured in customers per minute.

On completion of service at N$_1$,
50\% of the jobs queue again at N$_1$,
25\% of the jobs queue at N$_2$, and
25\% of the jobs queue at N$_3$.

On completion of service at N$_2$,
25\% of the jobs queue at N$_1$ and
75\% of the jobs queue at N$_3$.

On completion of service at N$_3$,
50\% of the jobs queue again at N$_3$
and 50\% of the jobs leave the system.

\begin{enumerate}

\item Show that the effective flow rates at the three nodes are:
\hfill {\bf [2 marks]}\\

\begin{center}
$\Lambda_1$ = 2 \hspace*{0.5in} 
$\Lambda_2$ = 1 \hspace*{0.5in} 
$\Lambda_3$ = 4
\end{center}

measured in customers per minute.

\item Hence find individual node utilisations   \hfill {\bf [2 marks]}\\

\item ... the expected number of customers at each node, and 
the expected number of customers in the system as a whole,
\hfill {\bf [2 marks]}\\


\item ... and average time spent by a customer in the system.  
\hfill {\bf [2 marks]}\\

\end{enumerate}

\newpage

\item Consider four variants of the same theme. \hfill {\bf [6 marks]}\\

\begin{enumerate}

\item M/M/1 system (one queue, one server with
          mean service rate of 1 job/minute)

\item an M/M/1/2 system (at most 2 jobs in the system,
so the maximum queue length is 1,
one server with mean service rate of 1 job/minute)

\item M/M/2 system (one queue, two servers each with a
          mean service rate of 0.5 jobs/minute)

\item M/M/2/2 system (one queue of maximum capacity 1, 
    two servers each with a mean service rate of 0.5 jobs/minute)

\end{enumerate}

For arrival rates $\lambda$ of 0.1, 0.5, and 0.9 jobs per minute,
tabulate $\rho$, N, and T (the utilisation, 
expected number in the system, and
expected through times).
Any conclusions you can draw?


The formulae are given in your main notes for MM1 and MMc;
as a prelim for this coursework for MM1/2 and you just did MM2/2.
Here they are:

\begin{center}
\begin{tabular}{lc lc lc l}
{\bf MM1} & & $\rho$ = ${\lambda}$/${\mu}$ &&
          N  = $\rho$/(1-$\rho$)&&
          T =  E[s]/(1-$\rho$)\\[0.1in]

{\bf MM1/2} & & $\rho$ = ${\lambda}$/${\mu}$ &&
          N  = $\rho$(1+2$\rho$)/(1+$\rho$+$\rho^2$)&&
          T =  E[s](1+2$\rho$)/(1+$\rho$)\\[0.1in]

{\bf MM2} & & $\rho$ = ${\lambda}$/2${\mu}$ &&
          N  = 2$\rho$/(1-$\rho^2$)&&
          T =  E[s]/(1-$\rho^2$)\\[0.1in]

{\bf MM2/2} && $\rho$ = $\lambda$/2$\mu$&&
N     =  2$\rho$(1+2$\rho$)/(1+2$\rho$+2$\rho^2$) &&
T     =  E[s] \\[0.1in]
\end{tabular}
\end{center}

{\bf Don't forget that $\rho$ = $\lambda$/2$\mu$ (or $\lambda$E[s]/2)
for both MM2 and MM2/2}. 

Here are the results:


\small
\begin{verbatim}
MM1

lam    mu     rho    Es       N           T
0.1    1.0    0.1    1.0    0.1111      1.1111
0.5    1.0    0.5    1.0    1.0000      2.0000
0.9    1.0    0.9    1.0    9.0000     10.0000

MM1/2

lam    mu     rho    Es       N           T
0.1    1.0    0.1    1.0    0.1081      1.0909
0.5    1.0    0.5    1.0    0.5714      1.3333
0.9    1.0    0.9    1.0    0.9299      1.4737

MM2

lam    mu     rho    Es       N           T
0.1    0.5    0.1    2.0    0.2020      2.0202
0.5    0.5    0.5    2.0    1.3333      2.6667
0.9    0.5    0.9    2.0    9.4737     10.5263


MM22

lam    mu     rho    Es       N           T
0.1    0.5    0.1    2.0    0.1967       2.0
0.5    0.5    0.5    2.0    0.8000       2.0
0.9    0.5    0.9    2.0    1.1403       2.0
\end{verbatim}
\normalsize

Conclusions: 

MM1 is better than MM2 with half speed servers (aside:
there will always be a need for fast cpu's).
This carries over to MM1/2 better than MM2/2 with half speed servers.
MMc (infinite buffer) is a poor approximation to MMc/n (finite buffer).

Here is my code (in SML with l for $\lambda$, m for $\mu$, and r for $\rho$):

\small
\begin{verbatim}
fun mm1 l m = let val r = l/m
                  val Es = 1/m
                  val N = r/(1-r)
                  val T = Es/(1-r)
              in
                  (l, m, r, Es, N, T)
              end;

fun mm12 l m = let val r = l/m
                  val Es = 1/m
                  val N = r*(1+2*r)/(1+r+r*r)
                  val T = Es*(1+2*r)/(1+r)
              in
                  (l, m, r, Es, N, T)
              end;

fun mm2 l m = let val r = l/(2*m)
                  val Es = 1/m
                  val N = 2*r/(1-r*r)
                  val T = Es/(1-r*r)
              in
                  (l, m, r, Es, N, T)
              end;

fun mm22 l m = let val r = l/(2*m)
                  val Es = 1/m
                  val N = 2*r*(1+2*r)/(1+2*r+2*r*r)
                  val T = Es
               in
                   (l, m, r, Es, N, T)
               end;
\end{verbatim}
\normalsize

Typical calls: mm1 0.1 1.0; mm12 0.1 1.0; mm2 0.1 0.5; mm22 0.1 0.5;


\newpage


\item Consider an open Jackson network comprising
three M/M/1 systems, N$_1$, N$_2$, N$_3$, with mean service rates 

\begin{center}
$\mu_1$ = 4 \hspace*{0.5in} $\mu_2$ = 4 \hspace*{0.5in} $\mu_3$ = 6
\end{center}

and mean customer arrival rates of

\begin{center}
$\lambda_1$ = 0.75 \hspace*{0.5in} 
$\lambda_2$ = 0.50 \hspace*{0.5in} 
$\lambda_3$ = 0.75
\end{center}

both measured in customers per minute.

On completion of service at N$_1$,
50\% of the jobs queue again at N$_1$,
25\% of the jobs queue at N$_2$, and
25\% of the jobs queue at N$_3$.

On completion of service at N$_2$,
25\% of the jobs queue at N$_1$ and
75\% of the jobs queue at N$_3$.

On completion of service at N$_3$,
50\% of the jobs queue again at N$_3$
and 50\% of the jobs leave the system.

\begin{enumerate}

\item Show that the effective flow rates at the three nodes are:
\hfill {\bf [2 marks]}\\

\begin{center}
$\Lambda_1$ = 2 \hspace*{0.5in} 
$\Lambda_2$ = 1 \hspace*{0.5in} 
$\Lambda_3$ = 4
\end{center}

measured in customers per minute.

{\sf 
The flow equations are :

\begin{tabular}{l c r c r c rc rc}
F1 & = & 0.75 & + & 0.50F1 & +  & 0.25F2 &  & \\
F2 & = & 0.50 & + & 0.25F1 &    &         &  & \\
F3 & = & 0.75 &   & 0.25F1 & +  & 0.75F2 & + & 0.50F3 \\
\end{tabular}

>From the equations for F1 and F2;

F1  =  0.75  +  0.50F1  +   0.25F2 and F2 = 0.50 + 0.25 F1.\\
Multiplying through by 4:\\
4F1  =  3 +  2F1  +   F2 and 4F2 = 2 + F1.\\
>From which 2F1 = 3 + F2 and   4F2 = 2 + F1.\\
>From which F1 = 2 and F2 = 1.\\
Substituting into the last equation for F3, F3 = 4.\\

Check: total flow in is 0.75 + 0.5 + 0.75 = 2 jobs/min.
Flow in = flow out = 0.5F3 = 2. Yes.

}


\item Hence find individual node utilisations   \hfill {\bf [2 marks]}\\

Utilisations are 

$rho_1$ = 0.5; $rho_2$ = 0.25; $rho_3$ = 0.66; 

\item ... the expected number of customers at each node, and 
the expected number of customers in the system as a whole,
\hfill {\bf [2 marks]}\\
N = $\rho$/(1-$\rho$):


$N_1$ = 1.00; $rho_2$ = 0.33; $rho_3$ = 2.00; 

N in system = 3.33


\item ... and average time spent by a customer in the system.  


By Little's Law, T = N/($\lambda_1$ + $\lambda_2$ + $\lambda_3$) =
3.33/2 = 1.67 minutes.


\hfill {\bf [2 marks]}\\

\end{enumerate}
\end{enumerate}


\begin{enumerate}

\item  Let S$_k$ signify the state of an {\em M/M/1} system when
there are {\em k} customers in the system,
and P$_k$ the probability of the system being in that state.
Clearly, P$_0$ + P$_1$ + P$_2$ + ... = 1.
Using the equations of job flow balance,
or otherwise, show that P$_k$ = $\rho^k$P$_0$,
and hence that P$_0$ = 1-$\rho$, where $\rho$ is the utilisation.

Hence or otherwise show that in an M/M/1 system,
the expected number of customers N is  $\rho$/(1-$\rho$) and 
the expected value of the customer through-time T is
E[s]/(1-$\rho$), where E[s] is the expected service time.


\item  Consider a (Jackson) network consisting of three 
interconnected M/M/1 systems.

All external jobs entering the system do so at node 1.
Their arrival rate is 10 jobs/minute.

\begin{itemize}
\item from node 1, jobs go to node 2 with probability 0.4 
and to node 3 with probability 0.4,
and leave the network with probability 0.2 

The mean service rate at node 1, $\mu_1$, is 100 jobs/minute

\item from node 2, jobs go to node 1 with probability 1.0.
The mean service rate at node 2, $\mu_2$, is 80 jobs/minute


\item from node 3, jobs go to node 1 with probability 1.0
The mean service rate at node 3, $\mu_2$, is 25 jobs/minute
\end{itemize}


\begin{enumerate}
\item Draw a picture of the network showing the three nodes.

\item Write down the flow rate equations for each node,
and show that 
\begin{center}
$\Lambda_1$ = 50 jobs/minute; $\Lambda_2$ =  $\Lambda_3$ = 20 jobs/minute
\end{center}

where $\Lambda_1$, $\Lambda_2$, and $\Lambda_3$
are the actual flow rates at the 4 nodes respectively

\item What are the utilisations of the three nodes?

\item Which node is the bottleneck in this system?
At what flow rate of external jobs would the system saturate?

\end{enumerate}



\item Queueing

\begin{enumerate}

\item M/M/1

Job flow balance gives us that
P$_1$ = $\rho$.P$_0$;  P$_2$ = $\rho$.P$_1$;  P$_3$ = $\rho$.P$_2$;
...\\
Hence P$_k$ = $\rho^k$P$_0$.
And P$_0$ = 1/(1-$\rho$)


 N = 0.P$_0$ + 1.P$_1$ + 2.P$_2$ + ...
= $\rho$.(1 + 2.$\rho$  + 3.$\rho^2$ + ...).P$_0$
= $\rho$.P$_0$/(1-$\rho$)$^2$
= $\rho$/(1-$\rho$)

T follows from Little's Law.

\item Jackson
\begin{enumerate}
\item Draw a picture of the network showing the three nodes.

\item  $\Lambda_1$ = 10 + $\Lambda_2$ + $\Lambda_3$

$\Lambda_2$ = 0.4.$\Lambda_1$ 

$\Lambda_3$ = 0.4.$\Lambda_1$ 

Hence $\Lambda_1$ = 10 + 0.4.$\Lambda_1$ + 0.4.$\Lambda_1$

Or 0.2.$\Lambda_1$ = 10

Hence $\Lambda_1$ = 50 jobs/minute; 
$\Lambda_2$ =  $\Lambda_3$ = 20 jobs/minute

\item What are the utilisations of the three nodes?

$\rho_1$ = 50/100 = 0.50; $\rho_2$ = 20/80 = 0.25; $\rho_3$ = 20/25 = 0.80

\item The bottleneck device is Node$_3$ at 80\%.
The utilisations rise with the external flow-in until 
the bottleneck device has 100\% utilisation.
That's Node$_3$ when the external flow rates has risen from 10
jobs/min to 10*100/80 = 12.5 jobs/minute.

\end{enumerate}
\end{enumerate}


{\bf Queueing practice}


\begin{enumerate}

\item Given the M/M/2/2 system transitions:

\begin{figure}[hbt]\centering
\setlength{\unitlength}{0.16mm}
\begin{picture}(720,120)(0, 60)

\thicklines

\put(120,  120){\oval(60, 60)} 
\put(120,  120){\makebox(0, 0){\sf 0}}
\put(150,  140){\vector(1, 0){60}} 
\put(180,  150){\makebox(0, 0){$\lambda$}}


\put(240,  120){\oval(60, 60)} 
\put(240,  120){\makebox(0, 0){\sf 1}}
\put(270,  140){\vector(1, 0){60}} 
\put(300,  150){\makebox(0, 0){$\lambda$}}
\put(210,  100){\vector(-1, 0){60}} 
\put(180,   80){\makebox(0, 0){$\mu$}}

\put(360,  120){\oval(60, 60)} 
\put(360,  120){\makebox(0, 0){\sf 2}}
\put(330,  100){\vector(-1, 0){60}} 
\put(300,   80){\makebox(0, 0){2*$\mu$}}

\end{picture}
\label{M/M/1/K}
\end{figure}

\begin{enumerate}

\item show that P$_0$, the probability of no customers in the system =
1/(1 + 2$\rho$ + 2$\rho^2$), where {\bf $\rho$ = $\lambda$/2$\mu$},
and derive formulae for P$_1$ and P$_2$. 

\item show that N, the average number of customers in the
system, is given by N =  2$\rho$$\frac{1+2\rho}{1 + 2\rho + 2\rho^2}$


\item derive a formula for T, the average customer through time 


\end{enumerate}

\item Consider an open Jackson network comprising
three M/M/1 systems, N$_1$, N$_2$, N$_3$, with mean service rates 

\begin{center}
$\mu_1$ = 4 \hspace*{0.5in} $\mu_2$ = 4 \hspace*{0.5in} $\mu_3$ = 6
\end{center}

and mean customer arrival rates of

\begin{center}
$\lambda_1$ = 0.75 \hspace*{0.5in} 
$\lambda_2$ = 0.50 \hspace*{0.5in} 
$\lambda_3$ = 0.75
\end{center}

both measured in customers per minute.

On completion of service at N$_1$,
50\% of the jobs queue again at N$_1$,
25\% of the jobs queue at N$_2$, and
25\% of the jobs queue at N$_3$.

On completion of service at N$_2$,
25\% of the jobs queue at N$_1$ and
75\% of the jobs queue at N$_3$.

On completion of service at N$_3$,
50\% of the jobs queue again at N$_3$
and 50\% of the jobs leave the system.

\begin{enumerate}

\item Show that the effective flow rates at the three nodes are:
\hfill {\bf [3 marks]}\\

\begin{center}
$\Lambda_1$ = 2 \hspace*{0.5in} 
$\Lambda_2$ = 1 \hspace*{0.5in} 
$\Lambda_3$ = 4
\end{center}

measured in customers per minute.

\item Hence find individual node utilisations   

\item ... the expected number of customers at each node, and 
the expected number of customers in the system as a whole,



\item ... and average time spent by a customer in the system.  


\end{enumerate}
\end{enumerate}
\end{enumerate}

\end{document}

